% generated from papers/XXXI-rayclass/paper_XXXI_rayclass.tex -- edit the paper, then rerun assemble.py
\chapter[The ray-class repair for affine systems]{The ray-class repair for affine systems:\\ the class group as the index of definable scalings}\label{ch:XXXI}
\renewcommand{\Ns}{N}
\chapterorigin{This chapter is Paper XXXI of the series (\texttt{papers/XXXI-rayclass/}).}
\begin{summary}
Chapter~\ref{ch:XXX} computed the K-theory of the affine boundary quotient at a principal prime and
located the class group ``at the entrance'': for non-principal $\pp$ the affine monoid of
Chapter~\ref{ch:II} (its \S6) is not well defined over $\OK$. This paper performs the repair and
proves that the obstruction is itself the invariant. Let $h=\ord([\pp]\in\Cl_K)$ and
$\pp^h=(\pi)$. The globally definable scalings form exactly the subsemigroup
$h\N$ of the local valuation semigroup: a scaling of depth $j$ exists over $\OK$ if and
only if $h\mid j$ (Theorem 1.2). The minimal repaired system is therefore the $h$-step
system $C(\Op)\rtimes(\OK\rtimes_\pi\N)$, well defined for every $\pp$, reducing to Paper
XXX when $h=1$; its boundary quotient has dilated form
$C_0(K_\pp)\rtimes(\mathcal O_{K,\{\pp\}}\rtimes_\pi\Z)$, and the entire exterior-tower
Pimsner--Voiculescu computation of Chapter~\ref{ch:XXX} applies verbatim with $(\pi,q^h)$ in place of
$(\pi,q)$, $q=\Ns\pp$. Consequences: the class of the unit generates a torsion summand of
\emph{exact} order $q^h-1=\Ns\pp^{\,h}-1$ --- the order of $\pp$ in the class group
survives in $K_0$ as the exact torsion order of $[1]$, and the published divisibility
statement of Chapter~\ref{ch:II} must read $(\Ns\pp^{h}-1)[1]=0$ in the repaired system; the
intermediate layers contribute $\Z/|q^h\pm1-\Tr\pi|'$ for quadratic fields, computed
explicitly in $\Q(\sqrt{-5})$ at the non-principal prime above $3$ ($K_0\supseteq
\Z\oplus\Z/8$ with $[1]$ of order $8$, $K_1\supseteq\Z\oplus\Z/2$). What the K-groups
alone retain is the composite $q^h$; separating $q$ from $h$, the Galois descent from the
Hilbert class field (where capitulation principalizes $\pp$ and $\mathrm{Gal}\cong\Cl_K$
acts), and the several-prime case are recorded as the remaining problems.
\end{summary}


\section{The bottleneck, and the scaling-index theorem}\label{XXXI:sec:index}

Fix a prime $\pp$ of $K$, $q=\Ns\pp$, and let $h$ be the order of $[\pp]$ in the class
group $\Cl_K$, so that $\pp^h=(\pi)$ is principal and no smaller positive power is. The
affine monoid of Chapter~\ref{ch:II} (its \S6) acts by $x\mapsto\gamma x+b$ with $b\in\OK$ and
$\gamma$ a chosen uniformizing multiplier; the composition law requires $\gamma\OK\subseteq
\OK$ with $\gamma$ supported at $\pp$, i.e.\ $\gamma\in\OK$ with $(\gamma)=\pp^{\,j}$ for
some $j\ge1$.

\begin{lemma}\label{XXXI:lem:exist}
An element $\gamma\in\OK$ with $(\gamma)=\pp^{\,j}$ exists if and only if $h\mid j$.
\end{lemma}

\begin{proof}
$(\gamma)=\pp^{\,j}$ says $[\pp]^{\,j}$ is trivial in $\Cl_K$, which is $h\mid j$;
conversely $\pi^{j/h}$ works.
\end{proof}

\begin{theorem}[The class group is the index of definable scalings]\label{XXXI:thm:index}
The set of depths $j$ for which the affine system of Chapter~\ref{ch:II}'s \S6 admits a well-defined
scaling of depth $j$ over $\OK$ is exactly $h\N$. The minimal well-defined affine system
at $\pp$ is therefore the $h$-step system
\[
\Aaff_{K,\{\pp\}}:=C(\Op)\rtimes\bigl(\OK\rtimes_\pi\N\bigr),\qquad(\pi)=\pp^{h},
\]
with boundary quotient defined by the Cuntz relation over the $q^h$ translates
$\sum_{a\in\OK/\pp^{h}}\mu_a\mu_a^*=1$; it reduces to the system of Chapter~\ref{ch:XXX} when $h=1$.
Any choice of local uniformizer at intermediate depths $0<j<h$ carries the translation
module $\OK$ to a fractional module in the ideal class $[\pp]^{-j}\ne1$ and leaves the
category of $\OK$-affine systems.
\end{theorem}

\begin{proof}
Lemma~\ref{XXXI:lem:exist} gives the depths. For the last statement, a local uniformizer
$\varpi\in K_\pp$ of depth $j$ maps the dense image of a fractional ideal $\mathfrak a$ to
that of an ideal in the class $[\pp]^{\,j}[\mathfrak a]$ up to prime-to-$\pp$ global
rescalings, which act by homeomorphisms of $K_\pp$; the marked class is an invariant of
the pair (translation module, ambient $K_\pp$) modulo such rescalings, and it changes
unless $h\mid j$.
\end{proof}

\begin{remark}
Theorem~\ref{XXXI:thm:index} upgrades the observation of Chapter~\ref{ch:XXX} --- ``the class group is at
the entrance'' --- to a structure theorem: the entrance fee is exactly $h$. Route B of the
repair, passing to the Hilbert class field $H$ where the extended ideal $\pp\mathcal O_H$
is principal by the principal ideal theorem \cite[VI.7]{Neukirch} and
$\mathrm{Gal}(H/K)\cong\Cl_K$ acts, defines a one-step system upstairs whose Galois
descent should remember the group $\Cl_K$ itself rather than only $h$; its K-theory is
Problem (ii) of \S\ref{XXXI:sec:verdict}.
\end{remark}

\section{K-theory of the repaired boundary}\label{XXXI:sec:ktheory}

\begin{theorem}\label{XXXI:thm:ktheory}
Let $n=[K:\Q]$ and $\Ns_{K/\Q}\pi=\pm q^{h}$. The dilated form of the repaired boundary is
\[
\partial\Aaff_{K,\{\pp\}}\otimes\K\ \cong\
C_0(K_\pp)\rtimes\bigl(\mathcal O_{K,\{\pp\}}\rtimes_\pi\Z\bigr),
\]
where $\mathcal O_{K,\{\pp\}}=\bigcup_m\pi^{-m}\OK$ is the ring of $\{\pp\}$-integers,
dense in $K_\pp$; and the exterior-tower computation of Chapter~\ref{ch:XXX} applies verbatim with
$q^h$ in place of $q$: with $M_r=\varinjlim(\Lambda^r\OK,\,t_r)$,
$t_r=\pm(\Lambda^{n-r}\pi)^\vee$, $\alpha_*=q^{-h}\Lambda^r\pi$, and
$T_r=\operatorname{coker}(1-\alpha_*|M_r)$, one has, when $\Ns\pi=+q^h$ and
$\pi^A\ne q^h$ for all proper archimedean subsets $A$,
\[
K_0\cong\Z\oplus\bigoplus_{r\,\mathrm{even}}T_r,\qquad
K_1\cong\Z\oplus\bigoplus_{r\,\mathrm{odd}}T_r,\qquad
T_0=\Z/(q^{h}-1)\ \text{generated by }[1],
\]
(the two $\Z$'s replaced by a single $\Z/2$ in parity $n$ when $\Ns\pi=-q^h$), the orders
of the middle layers being the regularized $|\det(q^h\cdot1-\Lambda^r\pi)|$ as in Paper
XXX. In particular:
\begin{enumerate}
\item \emph{The order of $\pp$ in the class group survives in $K_0$}: the class $[1]$ has
exact order $\Ns\pp^{\,h}-1$. The divisibility statement of Chapter~\ref{ch:II}'s \S6 must
accordingly read $(\Ns\pp^{h}-1)[1]=0$ for non-principal $\pp$, and it is attained.
\item For $K$ quadratic, $K_1(\partial\Aaff)\supseteq\Z/\bigl|q^h\pm1-\Tr\pi\bigr|'$
(prime-to-$q$ part), the sign that of $\Ns\pi$.
\end{enumerate}
\end{theorem}

\begin{proof}
The dilation, the tower, the two calibrations and the Pimsner--Voiculescu \cite{PV} assembly are
those of Chapter~\ref{ch:XXX}, restated there in full; the only inputs were a global generator of the
scaling ideal and its norm, here $\pi$ and $\Ns(\pp^h)=q^h$. The exactness of the order of
$[1]$ is the $r=0$ layer: $1-\alpha_*=1-q^{-h}$ on $\Z[1/q]$ has cokernel $\Z/(q^h-1)$
generated by the class of $1_{\Op}$.
\end{proof}

\begin{example}[$\Q(\sqrt{-5})$ at the non-principal prime above $3$]\label{XXXI:ex:sqrt5}
$K=\Q(\sqrt{-5})$, $\Cl_K\cong\Z/2$, $\pp_3=(3,1+\sqrt{-5})$ non-principal, $h=2$,
$q=3$, and $\pp_3^{\,2}=(\pi)$ with $\pi=2-\sqrt{-5}$, $\Ns\pi=9=q^h$, $\Tr\pi=4$. Then:
$[1]$ has exact order $q^h-1=8$, so $K_0\supseteq\Z\oplus\Z/8$ with $[1]$ generating the
torsion; and the middle layer gives
$\Ns(q^h-\pi)=\Ns(7+\sqrt{-5})=54=2\cdot27$, of prime-to-$3$ part $2$, so
$K_1\supseteq\Z\oplus\Z/2$, consistently with
$|q^h+1-\Tr\pi|=|9+1-4|=6$ of prime-to-$3$ part $2$. Had $\pp_3$ been principal the
order of $[1]$ would have been $q-1=2$: the K-theory separates the two situations.
\end{example}

\begin{remark}[Separating $q$ from $h$: the middle layers generically do]\label{XXXI:rem:separate}
The order of $[1]$ sees only the composite $q^h$, and pairs realizing the same composite
exist: $\Q(\sqrt{-5})$ at $\pp_3$ ($q=3$, $h=2$) against $\Q(i)$ at the inert prime
$(3)$ ($q'=9$, $h'=1$, principal) both have $\ord[1]=8$. But the middle layers separate
them: for $\Q(i)$, $\pi=3$ acts as the scalar $\tfrac13$ on $M_1\cong\Z[1/3]^2$, so
$T_1=\operatorname{coker}(\tfrac23)^{\oplus2}\cong(\Z/2)^2$ of order $4$, against the
$\Z/2$ of Example~\ref{XXXI:ex:sqrt5}. So Problem (i) of \S\ref{XXXI:sec:verdict} is partially
resolved: the K-groups themselves generically separate $(q,h)$ through the trace data of
the intermediate exterior layers, and only pairs with matching characteristic data at
every layer remain for the dual coaction to distinguish.
\end{remark}

\begin{remark}[Unit-orbit sensitivity persists]\label{XXXI:rem:units}
As in Chapter~\ref{ch:XXX}, replacing $\pi$ by $u\pi$ for a unit $u$ of norm $1$ changes the middle
layers; for imaginary quadratic fields the unit group is finite and the ambiguity is a
root of unity, so the repaired invariant is canonical there up to finitely many choices
--- a small unexpected bonus of the imaginary quadratic case.
\end{remark}

\section{The arithmetic verdict, and what remains}\label{XXXI:sec:verdict}

The class group enters the K-theory as \emph{torsion, not rank}: the free ranks of
Theorem~\ref{XXXI:thm:ktheory} are those of Chapter~\ref{ch:XXX}, unchanged, while the order of the
distinguished class $[1]$ is promoted from $q-1$ to $q^{h}-1$. Three problems remain.
\emph{(i) Separating $q$ from $h$:} partially resolved by Remark~\ref{XXXI:rem:separate} ---
the middle layers generically separate; what remains is the exceptional locus of pairs
with matching layer data, where the dual coaction is the candidate. \emph{(ii) Galois descent from the Hilbert class field:}
capitulation principalizes $\pp$ upstairs with $\mathrm{Gal}(H/K)\cong\Cl_K$; the
equivariant K-theory of the descended system should remember $\Cl_K$ itself, and its
computation --- a genuine descent spectral sequence, not an iterated
Pimsner--Voiculescu --- is open. \emph{(iii) Several primes:} for $|S|\ge2$ the definable
scalings form the kernel of $\Z^S\to\Cl_K$, $e\mapsto\prod\pp^{e_\pp}$, a finite-index
subgroup whose index divides $|\Cl_K|$ and whose K-theoretic footprint, layer by layer, is
the natural continuation of this paper.

